% ============================================================
% 50_typography_semantics.tex — Schriftmakros für Box-Typen & semantische Inhaltsmakros
% ============================================================

% Explizite bp-Grössen (wie Analysis) für präzise, font-unabhängige Kontrolle
\newcommand{\ZSFfontChapter}{\sffamily\bfseries\fontsize{8.1bp}{9.1bp}\selectfont}
\newcommand{\ZSFfontSection}{\sffamily\bfseries\fontsize{8.1bp}{9.1bp}\selectfont}
\newcommand{\ZSFfontBoxTitle}{\sffamily\bfseries\fontsize{8.1bp}{9.1bp}\selectfont}
\newcommand{\ZSFfontTableBoxTitle}{\sffamily\bfseries\fontsize{7.3bp}{8.1bp}\selectfont}
% chktex 6
\newcommand{\ZSFfontNote}{\sffamily\fontsize{6.7bp}{7.5bp}\itshape\selectfont}
\newcommand{\ZSFfontStatementTitle}{\bfseries}
\newcommand{\ZSFfontProcedureStep}{\bfseries}
\newcommand{\ZSFfontFactLead}{\bfseries}
\newcommand{\ZSFfontDanger}{\sffamily\bfseries}
% Rechtsbuendiges Meta-/Kategorie-Tag im Titelbalken (vgl. \ZSFTitleTag in
% styles/60_boxes.tex). Etwas kleiner und kursiv, begleitet den Titel.
\newcommand{\ZSFfontTitleTag}{\sffamily\fontsize{6.4bp}{7.0bp}\itshape\selectfont}

% Tabellen-Hilfsmakros
\newcommand{\TableCellNote}[1]{{\ZSFfontNote #1}}
\newcommand{\TableTitleBreak}{\newline}
\newcommand{\TableTitle}[2][]{%
  \textbf{#2}%
  \if\relax\detokenize{#1}\relax\else\TableTitleBreak\TableCellNote{#1}\fi
}

% Semantische Tabellen-Fontgrössen
% (werden als optionales Argument an ZSFtable/ZSFtableFlat/ZSFtablePlain
%  übergeben — siehe styles/20_tables.tex).
% Default = \ZSFfontTable. Für lange, dichte Tabellen (viele Zeilen,
% knapper Inhalt) \ZSFfontTableDense verwenden.
\newcommand{\ZSFfontTable}{\normalsize}
\newcommand{\ZSFfontTableDense}{\footnotesize}

% Semantische Inhaltsmakros (konsistente Skalierung und einfache Anpassung)
%
% \ZSFkeyword{...} — hebt zentrale Fachbegriffe im Fliesstext hervor.
% Verwendung:
%   * nur in \begin{runintext}, \begin{defbox}, \begin{warnbox} und
%     aehnlichen Text-Boxen, NIEMALS in Formeln, \formulanote oder
%     Überschriften (\SubsectionBar, \StartChapter).
%   * sparsam einsetzen (Richtwert: 1-3 pro runintext, 4-6 im
%     Kapitel-Intro, insgesamt ca. 10 pro Kapitel) -- es sollen echte
%     Fachbegriffe (Gesetze, Effekte, Bauteile, zentrale Grössen)
%     sein, keine generischen Substantive oder Verben.
% Rendering ist bewusst schlicht als \textbf{...}; Änderungen zentral
% hier vornehmen (z.B. Farbe, Underline) -- nicht in Kapiteln.
\newcommand{\ZSFkeyword}[1]{\textbf{#1}}
\newcommand{\ZSFScriptRef}[1]{\hfill\textcolor{\ChapterBarTextColor!80}{\small (S.#1)}}

% \ZSFhl{...} — hebt die EINE prüfungskritischste Aussage pro Box hervor
% (fett + unterstrichen). Sparsam: hoechstens einmal pro Block. Für
% Fachbegriffe \ZSFkeyword, für Stolperfallen \ZSFdanger nutzen.
\newcommand{\ZSFhl}[1]{\underline{\textbf{#1}}}

% \ZSFref{label} — kompakter interner Querverweis, gerendert als "(→ 6.2)".
% Verwendung:
%   * wenn eine Stelle ein Verfahren/Gesetz nutzt, das in einem ANDEREN
%     Kapitel erklaert wird.
%   * NICHT für Delegationen im selben Kapitel und nicht für elementare
%     Standardoperationen — sparsam einsetzen.
% Ziel-Label via optionales Argument von \SubsectionBar/\StartChapter setzen:
%   \SubsectionBar[sec:induktion-lenz]{Lenzsche Regel}
% \ref* (hyperref) unterdrueckt den Link, Nummern bleiben automatisch korrekt.
% Farbe = Palette-Farbe des ZIELkapitels (Wegweiser-Prinzip): wird automatisch
% aus der Kapitelnummer der Referenz abgeleitet ("6.6" -> ChapterColor6).
% Bei Kapiteln jenseits der Palette wird derselbe Modulo-Lookup wie bei
% \StartChapter genutzt, damit gewrappte Kapitel und Verweise konsistent sind.
% Aufrecht + fett statt kursiv für maximale Lesbarkeit bei 6.7bp.
\makeatletter
\def\ZSF@refExtractFirst#1#2\@nil{#1}
\def\ZSF@refExtractChap#1.#2\@nil{#1}
% Bestimmt \ZSF@refTargetColor aus dem Label; Fallback auf die aktive
% Kapitelfarbe im ersten latexmk-Pass (Label noch undefiniert).
\newcommand{\ZSF@refComputeTargetColor}[1]{%
  \@ifundefined{r@#1}{%
    \edef\ZSF@refTargetColor{\chaptercolor}%
  }{%
    \edef\ZSF@refTargetNum{\expandafter\expandafter\expandafter
      \ZSF@refExtractFirst\csname r@#1\endcsname\@nil}%
    \edef\ZSF@refTargetChap{\expandafter\ZSF@refExtractChap\ZSF@refTargetNum.\@nil}%
    \ZSFSetPaletteColorNameFromNumber{\ZSF@refTargetChap}{\ZSF@refTargetColor}%
  }%
}
\newcommand{\ZSFref}[1]{%
  \begingroup
  \ZSF@refComputeTargetColor{#1}%
  \hyperref[#1]{{\ZSFfontNote\upshape\bfseries\color{\ZSF@refTargetColor}($\rightarrow$\,\ref*{#1})}}%
  \endgroup
}
\newcommand{\ZSFsectionref}[1]{%
  \begingroup
  \ZSF@refComputeTargetColor{#1}%
  \hyperref[#1]{{\upshape\bfseries\color{\ZSF@refTargetColor}\ref*{#1}}}%
  \endgroup
}
\makeatother

\newcommand{\ZSFconclusion}[1]{%
  \par\vspace{-\parskip}\vspace{\ZSFspaceS}%
  \noindent$\Rightarrow$ #1%
  \par\vspace{-\parskip}\vspace{\ZSFspaceS}%
}

% Tabellen-Styling (benötigt [table]{xcolor} und array)
\newcommand{\ZSFrowColor}{\rowcolor{\FormulaboxBackColor}}
\newcommand{\ZSFtableRuleColor}{black!70}
\newcommand{\SetZSFtableRuleColor}[1]{\renewcommand{\ZSFtableRuleColor}{#1}}
\newcommand{\ZSFtableRuleWidth}{0.5pt}
\newcommand{\SetZSFtableRuleWidth}[1]{\renewcommand{\ZSFtableRuleWidth}{#1}}
\newcommand{\ZSFbeginTableRuleStyle}{%
  \begingroup
  \arrayrulecolor{\ZSFtableRuleColor}%
  \setlength{\arrayrulewidth}{\ZSFtableRuleWidth}%
}
\newcommand{\ZSFendTableRuleStyle}{\endgroup}
